\section{Soporte definido para pasos de instalaciones (PEN-30)}
\label{sec:sd4-pen30}
Para los treinta pasos SF de PEN-27, con sesenta dispositivos de PEN-33 se selecciona el detalle C-AJ-8309, revisión 25 de junio de 2024, como referencia de conjunto de paso, con placa Roxtec de al menos 6 mm y marco SF soldado de fábrica en ambas caras del muro. El listado incluye SF-6x1; prescribe soporte de hormigón armado de al menos 114 mm, no un espesor supuesto de 165 mm, o bloques de hormigón clasificados; abertura de hasta 0,4032 m$^2$ y dimensión máxima 635 mm. La placa solapa al menos 89 mm alrededor de la abertura y usa fijaciones de 8 mm por 57 mm con arandelas, a no más de 127 mm entre sí y dentro de 25 mm de cada esquina. Se aplican lana mineral preformada y fijaciones del fabricante dentro del vano y cordón perimetral de sellante de al menos 6 mm. Una junta improvisada de espuma no satisface el detalle. Se exige marca y configuración del conjunto suministrado, sin afirmar que la selección equivale a una instalación ya certificada \parencite[pp.~1--4]{lote30-caj8309}.

Se fija como obra propia un soporte local por red A/B y sector: doce paneles de hormigón armado de 0,80 por 1,65 por 0,12 m, dos por sector sobre rutas independientes. Se distribuyen tres aberturas en el panel A y dos en B, por sector, con un marco por cara de cada abertura: datos/energía A, control, datos/energía B; el control no une físicamente alimentación ni datos de ambos caminos. Cada marco conserva abertura individual de 201 por 298 mm, 0,059898 m$^2$, y placa mínima de reserva 379 por 476 mm, sin convertir la reserva en referencia de placa Roxtec comprada. Tres placas verticales ocupan 1,428 m y dejan 0,222 m para extremos dentro de la reserva de altura; posiciones, distancia a borde y armadura se ajustan al detalle ejecutivo sin cortar barras. Las placas de lados A y B no comparten panel, fijación ni soporte. El panel sustituye localmente el bastidor de 80 mm: 120 mm de hormigón, dos placas de 6 mm y 350 mm de aislamiento suman 482 mm, dentro de la reserva de muro de 500 mm. Los encuentros y la continuidad de blindaje/aislamiento se calculan y verifican antes de ejecutar; el listado de paso no clasifica esas juntas ni el muro completo.

Como cálculo propio a 2.400 kg/m$^3$, cada panel contiene 380,16 kg y los doce suman 4.561,92 kg de hormigón antes de armadura, placas y anclajes. Las sesenta placas mínimas de acero aportarían $60(0,379)(0,476)(0,006)(7850)=509,8217$ kg adicionales si son independientes de la piel BPR; PEN-33 desarrolla el volumen neto de vanos y la corrección de dos caras. Se incorporan las cargas a cimentación/estructura nueva, sin imputarlas como ya soportadas por el suelo existente ni usar la presión media del gabinete de baterías como cálculo del muro. Dimensionamiento de armadura, anclajes, sismo, posiciones y placa comercial final siguen pendientes del proyecto civil. Esta especificación desarrolla una solución elegida, pero no declara aprobación del soporte compuesto ni cumplimiento de todas las cláusulas de RT-06.04.

El listado fija F de dos horas, mientras T depende del servicio: PVC/EPR de energía y comunicaciones de hasta doce pares tienen máximo una hora; comunicaciones de cien pares No.~24 AWG, media hora; ciertos cables de tres/cuatro conductores XLPE/PVC, cero horas. La fibra PVC de hasta 6 mm no se extiende a otro diámetro o cubierta. No se convierte F en EI ni se atribuye T de dos horas a todos los enlaces. La lista de empaquetado identifica tipo/material/diámetro de cada cable y su rating; se mantiene pendiente cotejo de enlaces reales, módulos/placas y requisito térmico aplicable. Piso técnico/adaptación, cables de movimiento marino y auxiliares conservan sus brechas. El diseño FM 676 no sustituye este alcance: tiene servicios y prestaciones más restringidos \parencite[pp.~3--4]{lote30-caj8309}\parencite{lote30-fm676}.

\section{Variante láser actual elegida para bancos (INC-30)}
\label{sec:sd4-inc30}
Se seleccionan dos Xtralis VESDA-E VEP-A00-1P-UL para F04/F05 y uno de reserva fría de bancos; INC-36 sustituye los cuatro VLF ordinarios anteriores por cuatro VEP adicionales, sin duplicar inventario ni los tres repuestos fríos vigentes. La ficha primaria 36106\_07 de julio de 2025 identifica la variante de un tubo, tecnología láser y aplicación Class I Division 2 grupos A/B/C/D. Se corrige así la omisión de esa declaración en la revisión previa de VEP; VLC/ECO(Ex) descontinuados no son compra propuesta. La selección declara alcance del fabricante, sin sustituir la clasificación del recinto ni convertir Division 2 en Zone 1 o aprobación del circuito completo. Certificado FM 29131, condiciones de instalación del área, confirmador independiente y aptitud de accesorios siguen pendientes de cotejo íntegro antes de autorizar compra/energización \parencite[pp.~1--2]{lote30-vep36106}.

Cada VEP mide 350 por 225 por 135 mm, pesa 4,4 kg y recibe 18--30 V DC. La ficha da máximos de 0,57 A en reposo y 0,59 A en alarma: a 24 V se reservan 27,36/28,32 W para los dos, sin usar consumo nominal como máximo ni energizar el repuesto simultáneamente. La cámara funciona a 0--38 grados C y el aire muestreado a -20/+60 grados C, sin condensación. El tubo de cada banco y su escape vuelven al mismo sector, sin descargar gas a TI ni a un detector ordinario. Los ocho puntos/14 m de proyecto por banco se recalculan en ASPIRE para esta variante, orificios, sensibilidad y transporte máximo exigido; límite de longitud de catálogo no valida esos parámetros. La hoja UL requiere modelar todas las redes y sensibilidad por orificio. La confirmación independiente, señales supervisadas hacia Cheetah, clasificación normal/falla/posaislamiento y conservación de alimentación siguen sin desarrollar completamente. RT-06.16 permanece pendiente \parencite[pp.~2--3]{lote30-vep36106}\parencite{lote30-vepinst}.

\section{Conjunto de recipientes y actuación reconciliado (EXT-30)}
\label{sec:sd4-ext30}
Se sustituye la mezcla precedente de recipientes DOT/EN por la cohorte EN del manual Fike 06-542 revisión 3, con Impulse y DFIA, que declara UL Ex4623/FM 3034180. Se selecciona para TI un 70-256 de 83 L con 55 kg; para bancos, dos 70-278 de 26 L con 24 kg cada uno; para UPS/circulación, tres 70-350 de 16 L con 11/11/15 kg. Se dispone idéntico conjunto de seis reemplazos cargados: doce recipientes, 140 kg activos y 140 kg de reserva. Son llenados de proyecto en incrementos de 0,5 kg, no órdenes realizadas; todos están dentro de intervalos del manual: 83 L 53,5--93; 26 L 17--29; 16 L 10,5--17,5 kg. El 70-264 DOT anterior queda sustituido y sus dimensiones no se trasladan al 70-350 \parencite[ficha~06-542-1-2, pp.~1--2]{lote30-fike542}.

Dimensiones/taras EN: 83 L diámetro 324 mm, altura 1.316 mm, tara 62 kg; 26 L diámetro 229 mm, altura 876,3 mm, tara 23 kg; 16 L diámetro 229 mm, altura 635 mm, tara 16 kg. Las taras de seis suman 156 kg; con 140 kg de agente, 296 kg por cohorte y 592 kg entre activos/reemplazos antes de gabinete, soportes y accesorios. Los seis pequeños pesan 170 kg con agente, frente a 162,2 kg de la variante DOT anterior. Las dos reservas de gabinete 3 por 0,8 por 1,8 m con frente de 1,20 m contienen seis celdas de 0,5 m cada una; se comprueban conexiones/maniobra y control 0--38 grados C. Los diámetros y alturas publicados caben geométricamente en las celdas, sin presentar las diferencias como holguras OEM. El conjunto mantiene instalación vertical con válvula hacia arriba \parencite[ficha~06-542-1-2, pp.~1--2]{lote30-fike542}.

Se seleccionan doce kits DFIA 70-311, seis activos y seis montados/disponibles en reemplazos, con 02-13279 e IVOS 02-14263, seis RCM 55-053 para los activos, seis interruptores de descarga 02-12534 y uno adicional de reserva por tipo RCM/interruptor. Se conecta un DFIA por RCM directamente a Cheetah Xi; no se intercala IRM por analogía con otra actuación. Cada descarga activa consume hasta 930 mA y supervisión 25,4 mA: seis requieren 5,58 A/0,1524 A a 24 V, 133,92/3,6576 W sólo de actuadores. Los reemplazos no se suman como seis disparos simultáneos adicionales. INC-36 selecciona cuatro FRPS para cuatro VEP ordinarios, cuatro lazos/124 direcciones y desarrolla balance individual del panel/SPS con módulos y avisadores; quedan curvas útiles, confirmadores de banco, rutas/tensión final y secuencia; la elección compatible no demuestra autonomía del panel \parencite[ficha~06-432-1-5-DF, pp.~1--6]{lote30-fike542}.

Todos los llenados tienen $m/V=2/3$ kg/m$^3$ y conservan la aproximación propia de concentración 8,0826 por ciento a 10 grados C y 8,9198 a 40 grados C. La concentración de diseño se justifica por combustible/norma y cálculo hidráulico, sin extender una comparación A/C a H$_2$ ni convertir la clasificación V-0 seleccionada de carcasa en aprobación de extinción. El manual publica redes preingenierizadas de una boquilla: para 16 L, 84-047-053 con tubo nominal 20 y máximo 13,7 m; para 26 L, 84-048-067 con nominal 25 y el mismo máximo, ocho codos incluidos y descuento por adicionales. Se identifican como referencias concretas, sin adoptarlas para trazados que no cumplan todas las condiciones, volumen/altura/orientación y llenado. Para el gabinete exterior y TI de 83 L se requiere red de ingeniería y cálculo Fike con orificios definitivos; no se promete boquilla XXXX. Retención, alivio, listado vigente del conjunto, hidráulica y secuencia segura gas/fuego mantienen RT-06.17 pendiente. GAS-28 impide sellar banco por simple corte AC o lectura instantánea baja \parencite[sec.~8, tabla~1]{lote30-fike542}.

